\documentclass{ximera}

\title{Average Velocity Activity}
\author{MATH 425: Calculus I}

\begin{document}
\begin{abstract}
    Working with the peers in your group, solve the following problems. Make sure to show and justify all your work. Make sure everyone in the group understands the solution and participates. Be prepared to report your answers to the whole class. 
\end{abstract}
\maketitle


\begin{exercise}
 If a ball is thrown upwards into the air with a velocity of 40 ft/s, its height in feet $t$ seconds later is given by $y=40t-16t^2$.  
  Find the average velocity for the time period beginning at 2 seconds and lasting 
     \begin{enumerate}
         \item 0.5 seconds
         \item 0.1 seconds
         \item 0.05 seconds
         \item 0.01 seconds
     \end{enumerate}
 Estimate the instantaneous velocity of the same ball when $t=2$ using your work above.  Explain what your answer says about what is happening to the ball at $t=2$. \\

 \textcolor{blue}{For 0.5 seconds: \\
     Interval: $[2, 2.5]$\\
     $y(2)=40(2)-16(2)^2=16$\\
     $y(2.5)=40(2.5)-16(2.5)^2=0$\\
     $m=\frac{y(2.5)-y(2)}{2.5-2}=\frac{0-16}{0.5}=-32$\\
     For 0.1 seconds: \\
     Interval: $[2, 2.1]$\\
     $y(2)=16$\\
     $y(2.1)=40(2.1)-16(2.1)^2=13.44$\\
     $m=\frac{y(2.1)-y(2)}{2.1-2}=\frac{13.44-16}{0.1}=-25.6$\\
     For 0.05 seconds: \\
     Interval: $[2, 2.05]$\\
     $y(2)=16$\\
     $y(2.05)=40(2.05)-16(2.05)^2=14.76$\\
     $m=\frac{y(2.05)-y(2)}{2.05-2}=\frac{14.76-16}{0.05}=-24.8$\\
     For 0.01 seconds: \\
     Interval: $[2, 2.01]$\\
     $y(2)=16$\\
     $y(2.01)=40(2.01)-16(2.01)^2=15.7584$\\
     $m=\frac{y(2.5)-y(2)}{2.5-2}=\frac{15.7584-16}{0.01}=-24.16$\\
     It appears that the slopes are approaching $m=-24$ for the slope of the tangent line.  This would mean that the ball is falling at 24ft/s.}
\end{exercise}

\begin{exercise}
Suppose that a frozen yogurt shop has the following price structure: 
 \begin{itemize}
     \item It costs \$1 for the bowl.
     \item Each ounce of yogurt costs \$0.50.
     \item Each ounce of toppings costs \$1.
 \end{itemize}
 Write a formula for the function, $C(x)$, that gives the cost in dollars of your bowl of frozen yogurt as a function of $x$, the ounces of yogurt you get, if you always choose to get half as much toppings as yogurt. What is the domain of this function in the context of the yogurt shop if the bowl can hold 24 ounces of yogurt?\\
 \textcolor{blue}{total cost=bowl+cost per ounce(yogurt)+cost per ounce(toppings)\\
 Ounces of toppings=1/2(ounces of yogurt)\\
 Let $y=$ ounces of yogurt.\\
 Total cost$=1+0.50y+1(.5y)=1+y$\\
 The unrestricted range of $T(y)=1+y$ is all real numbers, but we cannot get negative ounces of yogurt/toppings, so it makes sense to restrict it to be 0 or larger.  The cup also cannot hold infinitely much yogurt/toppings, so the total has to be less than or equal to 24.\\
 We can model that mathematically as $0\leq\text{ ounces of yogurt+ounces of toppings }\leq 24$. With our variables this becomes:
 $0\leq y+\frac12 y\leq 24$\\
 $0\leq \frac32 y\leq 24$\\
 $0\leq y\leq 16$\\
 The range is 0 to 16 ounces of yogurt. }
\end{exercise}

\end{document}